\section{Session Types via QTT}

\label{sec:sessions}

To illustrate how we can use quantities on a more
substantial example, let us consider how to use them to implement session
types. Session types~\cite{Honda1993,Honda2008} give types to communication channels, allowing us to
express exactly \emph{when} a message can be sent on a channel, ensuring that
communication protocols are implemented completely and correctly. There has
been extensive previous work on defining calculi for session types\footnote{A collection of implementations is available
at \url{http://groups.inf.ed.ac.uk/abcd/session-implementations.html}}. In
Idris 2, the combination of linear and dependent types means that we can
implement session types directly:

\begin{itemize}
\item \textbf{Linearity} means that a channel can only be accessed once, 
and once a message has been sent or received on a channel, the channel is
in a new state.
\item \textbf{Dependent Types} give us a way of describing protocols at the
type level, where progress on a channel can change according to values sent
on the channel.
\end{itemize}

A complete implementation of session types would be
a paper in itself, so we limit ourselves to dyadic session types in concurrent
communicating processes.  We assume that functions are \emph{total}, so
processes will not terminate early and communication will always
succeed. In a full library, dealing with \emph{distributed} as
well as \emph{concurrent} processes, we would also need to consider failures
such as timeouts and badly formed messages~\cite{Fowler2019}.

The key idea is to parameterise channels by the actions which will be
executed on the channel---that is, the messages which will be
sent and received---and to use channels linearly. We declare a
\texttt{Channel} type as follows:

\begin{code}
data Actions : Type where
     Send : (a : Type) -> (a -> Actions) -> Actions
     Recv : (a : Type) -> (a -> Actions) -> Actions
     Close : Actions

data Channel : Actions -> Type
\end{code}

Internally, \texttt{Channel} contains a message queue for bidirectional
communication.
Listing \ref{sessioniface} shows the types of functions for
initiating sessions, and sending and receiving messages.
In the type of \texttt{send}, we see that to send a value of type
\texttt{ty} we must have a channel in the state \texttt{Send ty next},
where \texttt{next} is a function that computes the rest of the protocol.
The type of \texttt{recv} shows that we compute the rest of the protocol
by inspecting the value received. We initiate concurrent sessions with
\texttt{fork}, and will discuss the details of this shortly.

\begin{code}[float=h,frame=single,caption={Initiating and executing
concurrent sessions},label=sessioniface]
send : (1 chan : Channel (Send ty next)) -> (val : ty) ->
       L {use=1} (Channel (next val))
recv : (1 chan : Channel (Recv ty next)) ->
       L {use=1} (Res ty (\val => Channel (next val)))
close : (1 chan : Channel Close) -> L ()
fork : ((1 chan : Server p) -> L ()) -> L {use=1} (Client p)
\end{code}

First, let us see how to describe dyadic protocols such that a \emph{client}
and \emph{server} are guaranteed to be synchronised. We describe protocols
via a \emph{global} session type:

\begin{code}
data Protocol : Type -> Type where
     Request : (a : Type) -> Protocol a
     Respond : (a : Type) -> Protocol a
     (>>=) : Protocol a -> (a -> Protocol b) -> Protocol b
     Done : Protocol ()
\end{code}

A protocol involves a sequence of \texttt{Request}s from a client to
a server, and \texttt{Response}s from the server back to the client.
For example, we could define a protocol (Listing~\ref{utilsproto})
in which a client sends a \texttt{Command} to either \texttt{Add} a pair of
\texttt{Int}s or \texttt{Reverse} a \texttt{String}.

\begin{code}[float=h,frame=single,caption={A global session type describing
a protocol where a client can request either adding two \texttt{Int}s or
reversing a \texttt{String}},label=utilsproto]
data Command = Add | Reverse

Utils : Protocol ()
Utils = do cmd <- Request Command
           case cmd of
                Add => do Request (Int, Int)
                          Respond Int
                          Done
                Reverse => do Request String
                              Respond String
                              Done
\end{code}

\texttt{Protocol} is a DSL for describing communication
patterns. Embedding it in a dependently typed host language gives us
dependent session types for free. We use the embedding to our advantage 
in a small way, by having the protocol depend on \texttt{cmd}, the command
sent by the client.
%
We can write functions to calculate the protocol for the client and the
server: 

\begin{code}
AsClient, AsServer : Protocol a -> Actions
\end{code}

We omit the definitions, but each translates \texttt{Request} and
\texttt{Response} directly to the appropriate \texttt{Send} or
\texttt{Receive} action.
We can see how \texttt{Utils} translates into a type for the client side by
running \texttt{AsClient Utils}:

\begin{code}
Send Command (\res => (ClientK 
    (case res of 
          Add => Request (Int, Int) >>= \_ => 
                 Respond Int >>= \_ Done
          Reverse => Request String >>= \_ => 
                     Respond String >>= \_ Done)
\end{code}

Most importantly, this shows us that the first client side operation must be
to send a \texttt{Command}. The rest of the type is calculated from the
command which is sent; \texttt{ClientK} is internal to \texttt{AsClient} and
calculates the continuation of the type. Using these, we can
define the type for \texttt{fork}.

\begin{code}
Client, Server : Protocol a -> Type
Client p = Channel (AsClient p)
Server p = Channel (AsServer p)

fork : ((1 chan : Server p) -> L ()) -> L {use=1} (Client p)
\end{code}

The type of \texttt{fork} ensures that the client and the server are working
to the same protocol, by calculating the channel type of each from the same
global protocol. Since each channel is linear, both ends of the protocol
must be run to completion.

\begin{code}[float=h,frame=single,caption={An implementation of a server
for the \texttt{Utils} protocol},label=utilserver]
utilServer : (1 chan : Server Utils) -> L ()
utilServer chan
    = do cmd # chan <- recv chan
         case cmd of
              Add => do (x, y) # chan <- recv chan
                        chan <- send chan (x + y)
                        close chan
              Reverse => do str # chan <- recv chan
                            chan <- send chan (reverse str)
                            close chan
\end{code}

Listing~\ref{utilserver} shows a complete implementation of a server for
the \texttt{Utils} protocol. However, we do not typically write a complete
implementation in one go. Idris 2's support for \emph{holes} means that it
is more convenient to write the server incrementally, in a type-driven
way. We begin with just a skeleton definition, and look at the hole for the
right hand side:

\begin{code}
utilServer : (1 chan : Server Utils) -> L ()
utilServer chan = ?utilServer_rhs

 1 chan : Channel (Recv Command (\res => ... ))
-------------------------------------
utilServer_rhs : L ()
\end{code}

Therefore, the first action on \texttt{chan} must be to receive a
\texttt{Command}:

\begin{code}
utilServer : (1 chan : Server Utils) -> L ()
utilServer chan
    = do cmd # chan <- recv chan
         ?utilServer_rhs

   cmd : Command
 1 chan : Channel (ServerK (case cmd of ...) ...)
------------------------------
utilServer_rhs : L ()                                       
\end{code}

We elide the full details of the type of \texttt{chan} at this stage, but at
the top level it suggests that we can make progress by a \texttt{case} split on
\texttt{cmd}:

\begin{code}
utilServer : (1 chan : Server Utils) -> L ()
utilServer chan
    = do cmd # chan <- recv chan
         case cmd of
              Add => ?process_add
              Reverse => ?process_reverse
\end{code}

We make essential use of dependent \texttt{case} here, in that both
branches have a different type which is computed from the value of the
scrutinee \texttt{cmd}, similarly to \FN{PrintfType} in
Section~\ref{sec:printf}.
Now, for each of the holes \texttt{process\_add} and \texttt{process\_reverse}
we see more concretely how the protocol should proceed. e.g. for
\texttt{process\_add}, we see we have to receive a pair of \texttt{Int}s, then
send an \texttt{Int}:

\begin{code}
   1 chan : Channel (Recv (Int, Int) (\res => 
                    (Send Int (\res => Close))))
     cmd : Command
  -------------------------------------
  process_add : L () 
\end{code}

Developing the server in this way shows programming as a \emph{dialogue}
with the type checker. Rather than trying to work out the complete program,
with increasing frustration as the type checker rejects our attempts, we
write the program step by step, and ask the type checker for more information
on the variables in scope and the required result.

\subsubsection*{Extension: Sending Channels over Channels}

A useful extension is to allow a server to start up more
worker processes to handle client requests. This would require sending the
server's \TC{Channel} endpoint to the worker process. However, we cannot do
this with \TC{send} as it stands, because the value sent must be of
multiplicity $\omega$, and the \TC{Channel} is linear. One way to support
this would be to refine \TC{Protocol} to allow flagging messages as linear,
then add:

\begin{code}
send1 : (1 chan : Channel (Send1 ty next)) -> (1 val : ty) ->
        L {use=1} (Channel (next val))
\end{code}

This takes advantage of QTT's ability to parameterise types by linear variables
like \texttt{val} here. A worker protocol, using the \FN{Utils} protocol
above, could then be described as follows, where a server forks a new 
worker process and immediately sends it the communication \TC{Channel} for
the client:

\begin{code}
MakeWorker : Protocol ()
MakeWorker = do Request1 (Server Utils); Done
\end{code}

We leave full details of this implementation for future work. It is,
nevertheless, a minor adaptation of the session types library.
